\section{SF-E817: an actual singular balanced character and the
universal point-defect
obstruction}\label{SF-E817-sf-e817-an-actual-singular-balanced-character-and-the-universal-point-defect-obstruction}

Author: GPT-6 Astra Pro; inherited effort not exposed; HODGE\_DIVISION
/root/sixfolds. Date: 2026-11-21.

\textbf{Status:} Whole-proof mathematical PASS received from /root and
/root/imprimitive on2026-11-21 (GPT-6 Astra Pro; inherited effort not
exposed). Root read the complete written proof and exact certificate and
freshly reran that certificate. Imprimitive read the full proof, freshly
reconstructed the exterior calculation with independent tuple-wedge
code, and freshly reread the supplied S-20/S-21 categorical passages.
The author ran the displayed exact certificate separately. The theorem
concerns torsion-free sheaves with locally free double dual and
zero-dimensional quotient. It does not treat positive-dimensional
defects or singular reflexive hulls.

\subsection{1. Statement}\label{SF-E817-statement}

Let \(X\) be a polarized complex abelian fourfold with a genuine real
quadratic action, of holomorphic type \((2,2)\). Let
\(\theta_+,\theta_-\) be the components of its compatible polarization
\(\theta\), and put

\[
 \Delta=\int_X\frac{\theta^4}{24}>0,
 \qquad z_\pm=\frac{\theta_\pm^2}{2}.
 \tag{E817.1}
\]

Thus \(\int_Xz_+z_-=\Delta\). Let \(E\) be a torsion-free coherent sheaf
of positive rank \(r\), with locally free double dual
\(H=E^{\vee\vee}\), and suppose

\[
 0\longrightarrow E\longrightarrow H\longrightarrow T\longrightarrow0
 \tag{E817.2}
\]

has nonzero zero-dimensional quotient \(T\). Suppose that for a rational
algebraic divisor \(B\) and nonzero real rational \(q\),

\[
 e^B\operatorname{ch}(E)
      =r v_q,
 \qquad
 v_q=(1-qz_+)(1-qz_-).
 \tag{E817.3}
\]

Let \(G\subset X\times\widehat X\) be any effective finite subgroup with
a specified coherent translation/tensor action, and suppose that \(E\)
has an actual equivariant structure. Let \(s\ge1\) be the number of
\(G\)-orbits in the support of \(T\).

Then its full invariant degree-two trace has

\[
 \dim\ker\sigma^2_{E,G}\ge1+3s\ge4.
 \tag{E817.4}
\]

In particular \(E\) is not weakly \(G\)-semiregular. There is no
simplicity, stability, or polystability hypothesis on \(E\) or on \(H\).
Nonreduced finite-length quotients are allowed. The proof uses the
actual canonical double-dual sequence, not an assertion that its
subobject or quotient is a direct summand.

\subsection{2. The equivariant structures on the canonical
sequence}\label{SF-E817-the-equivariant-structures-on-the-canonical-sequence}

Translation and tensoring by a line bundle commute naturally with double
dual. The given coherent isomorphisms on \(E\) therefore induce coherent
isomorphisms on \(H=E^{\vee\vee}\), and the canonical map \(E\to H\) is
equivariant. Its quotient \(T\) inherits the actual coherent structure.
Associativity is preserved by these functorial operations; no new
projective linearization of an individual factor is introduced.

Every finite-length sheaf has a canonical decomposition into its pieces
at the finitely many support points, by the decomposition of its
Artinian support algebra. Grouping these pieces by translation orbits
gives actual equivariant summands

\[
 T=\bigoplus_{\nu=1}^{s}T_\nu.
 \tag{E817.5}
\]

Tensoring by a degree-zero line bundle does not change support. Distinct
support-orbit summands have disjoint support, hence all their cross Ext
groups vanish. Each \(T_\nu\) has positive length \(\ell_\nu\) and
character \(\ell_\nu[\mathrm{pt}]\).

\subsection{3. A canonical exact sequence on middle
self-Ext}\label{SF-E817-a-canonical-exact-sequence-on-middle-self-ext}

We first record two elementary vanishings on a smooth fourfold:

\[
 \operatorname{Ext}^j_X(H,T)=0\quad(j>0),
 \qquad
 \operatorname{Ext}^j_X(T,H)=0\quad(j<4).
 \tag{E817.6}
\]

For the first, \(H\) is locally free and \(H^\vee\otimes T\) has finite
support, so it has no higher cohomology. For the second, the local
Koszul resolution of a point in a regular local ring of dimension four
gives \(\operatorname{Ext}^j(k_p,H)=0\) for \(j<4\). A finite-length
sheaf has a finite filtration with such residue-field factors. The Ext
long exact sequence proves the same vanishing for every factor
extension. The local-to-global passage has no higher-support cohomology,
giving (E817.6). Negative Ext vanishes for sheaves as usual.
Finite-group invariants over \(\mathbf C\) are exact, so these
vanishings also hold for the specified equivariant Ext groups.

Use the actual equivariant triangle

\[
 T[-1]\xrightarrow{i}E\xrightarrow{p}H
       \xrightarrow{a}T.
 \tag{E817.7}
\]

Apply \(\operatorname{RHom}_G(H,-)\). The exact segment for \(j=2,3\) is

\[
 \operatorname{Ext}^{j-1}_G(H,T)
 \longrightarrow\operatorname{Ext}^{j}_G(H,E)
 \xrightarrow{p_*}\operatorname{Ext}^{j}_G(H,H)
 \longrightarrow\operatorname{Ext}^{j}_G(H,T).
\]

The exterior groups vanish, so

\[
 \operatorname{Ext}^{j}_G(H,E)
       \simeq\operatorname{Ext}^{j}_G(H,H)
       \qquad(j=2,3).
 \tag{E817.8}
\]

Next apply \(\operatorname{RHom}_G(T[-1],-)\). Its adjacent groups with
target \(H\) are \(\operatorname{Ext}^{j}_G(T,H)\) and
\(\operatorname{Ext}^{j+1}_G(T,H)\). For \(j=1,2\), both vanish by
(E817.6). Therefore

\[
 i_*:\operatorname{Ext}^{j}_G(T,T)
      \xrightarrow{\sim}\operatorname{Ext}^{j}_G(T[-1],E)
        \qquad(j=1,2).
 \tag{E817.9}
\]

Finally apply \(\operatorname{RHom}_G(-,E)\). The relevant complete
segment is

\[
 \begin{split}
 \operatorname{Ext}^1_G(T[-1],E)
 &\xrightarrow{\partial_1}\operatorname{Ext}^2_G(H,E)
 \longrightarrow\operatorname{Ext}^2_G(E,E)\\
 &\longrightarrow\operatorname{Ext}^2_G(T[-1],E)
 \xrightarrow{\partial_2}\operatorname{Ext}^3_G(H,E).
 \end{split}
 \tag{E817.10}
\]

The two connecting maps are zero. Indeed, using (E817.9), a degree-\(j\)
element, for \(j=1,2\), is represented by

\[
 f=i[j]\circ b[-1],\qquad b:T\longrightarrow T[j].
\]

Its connecting class is, up to sign,

\[
 \partial_j(f)=f[1]\circ a
             =i[j+1]\circ b\circ a.
 \tag{E817.11}
\]

The intermediate composite \(b\circ a\) belongs to
\(\operatorname{Ext}^j_G(H,T)=0\), by (E817.6). These are actual
equivariant morphisms, so the vanishing holds before any dimension
estimate.

Equations (E817.8)--(E817.11) yield the short exact sequence

\[
 0\longrightarrow\operatorname{Ext}^2_G(H,H)
   \longrightarrow\operatorname{Ext}^2_G(E,E)
   \longrightarrow\operatorname{Ext}^2_G(T,T)
   \longrightarrow0.
 \tag{E817.12}
\]

The sequence is not claimed to split canonically. It gives the exact
dimension identity

\[
 e_2^G(E)=e_2^G(H)+e_2^G(T).
 \tag{E817.13}
\]

No degeneration of a spectral sequence, formality assertion, or
stability condition is used.

\subsection{4. The complete cap-rank calculation after adding a point
class}\label{SF-E817-the-complete-cap-rank-calculation-after-adding-a-point-class}

Let \(\ell=\operatorname{length}(T)>0\). Since an exponential fixes the
top point class, (E817.2)--(E817.3) give

\[
 e^B\operatorname{ch}(H)
   =r(1-qz_+)(1-qz_-)+\frac{\ell}{\Delta}z_+z_-.
 \tag{E817.14}
\]

We compute its full degree-two Hochschild cap rank. Let

\[
 W=W_+\oplus W_-,\qquad
 W_\pm=L_\pm\oplus M_\pm,
 \qquad\dim L_\pm=\dim M_\pm=2,
\]

where \(L_\pm\) and \(M_\pm\) are the translation and tensor planes for
the respective real embeddings. Choose nonzero volume bivectors
\(\alpha_\pm\in\bigwedge^2L_\pm\) and \(\beta_\pm\in\bigwedge^2M_\pm\)
so that the Fock representatives of the unit and of \(z_\pm\) are
\(\alpha_\pm\) and \(\beta_\pm\). The full character in (E817.14)
becomes

\[
 w=a\alpha_+\alpha_-+b\beta_+\alpha_-
       +c\alpha_+\beta_-+d\beta_+\beta_-,
 \tag{E817.15}
\]

with

\[
 a=r,\qquad b=c=-rq,\qquad
 d=rq^2+\frac{\ell}{\Delta},\qquad
 ad-bc=\frac{r\ell}{\Delta}\ne0.
 \tag{E817.16}
\]

All four coefficients are nonzero. The multiplication map
\(\bigwedge^2W\to\bigwedge^6W\), \(u\mapsto u\wedge w\), has the
following three independent bidegree blocks.

On \(\bigwedge^2W_+\), only the two directions \(\alpha_+,\beta_+\)
contribute. Their images, in \(\bigwedge^4W_+\otimes\bigwedge^2W_-\),
have coefficient matrix

\[
 \begin{pmatrix}b&a\\d&c\end{pmatrix}.
\]

Its determinant is \(bc-ad\ne0\), so this block has rank two. The
analogous block on \(\bigwedge^2W_-\) has rank two as well.

On \(W_+\otimes W_-\), decompose each factor into its two
translation/tensor planes. The four resulting four-dimensional summands
map isomorphically into distinct summands of
\(\bigwedge^3W_+\otimes\bigwedge^3W_-\), multiplied respectively by
\(a,b,c,d\), up to signs. Hence the mixed block has rank sixteen. Thus

\[
 \operatorname{rank}c^2_{\operatorname{ch}(H)}=2+2+16=20.
 \tag{E817.17}
\]

The common divisor shear does not change this rank. When \(\ell=0\), the
two pure blocks instead have rank one each, recovering the balanced rank
eighteen. For a nonzero point character alone, only \(d\) is nonzero;
the same block calculation gives ranks one, one, and four, totaling six.

The calculation is universal in the coefficients. A nonzero
twenty-by-twenty minor can be chosen with determinant, up to sign,

\[
 (ad-bc)^2(abcd)^4.
 \tag{E817.18}
\]

An exact exterior script checks this generic minor and the stated ranks.
The proof of the universal rank is the three-block argument above.

\subsection{5. The invariant kernel lower
bound}\label{SF-E817-the-invariant-kernel-lower-bound}

For every actual equivariant perfect object, evaluation of Hochschild
degree two takes values in invariant Ext, because translation and
degree-zero tensor operations act trivially on Hochschild cohomology.
Evaluation followed by full trace is cap with its Chern character, up to
the fixed nonzero identity-sector normalization of supplied DEG\_L S-21.
Equation (E817.17) therefore gives

\[
 e_2^G(H)\ge20.
 \tag{E817.19}
\]

For each actual equivariant support-orbit summand \(T_\nu\), its
positive point character has cap rank six. Thus \(e_2^G(T_\nu)\ge6\).
The disjoint-support decomposition (E817.5) gives

\[
 e_2^G(T)\ge6s.
 \tag{E817.20}
\]

Combining these inequalities with (E817.13) yields

\[
 e:=e_2^G(E)\ge20+6s.
 \tag{E817.21}
\]

To obtain an actual trace-kernel bound without assuming semiregularity,
let \(h=\operatorname{rank}\sigma^2_{E,G}\) and \(k=e-h\). CY4
trace/evaluation adjunction says that the evaluation map
\(\rho^2_{E,G}\) also has rank \(h\). The Serre pairing on invariant Ext
is perfect, since the finite translation/tensor action preserves it and
averaging is exact. Cap with the balanced character has rank eighteen.
The elementary rank inequality for the composition gives

\[
 18=\operatorname{rank}(\sigma^2_{E,G}\rho^2_{E,G})
       \ge 2h-e=e-2k.
 \tag{E817.22}
\]

Equations (E817.21)--(E817.22) imply

\[
 k\ge\frac{e-18}{2}\ge1+3s,
\]

which is (E817.4). The bound remains positive for every effective finite
coherent action, including actions that correlate different support
points and carry nontrivial multiplicity representations.

\subsection{6. An explicit actual singular balanced
sheaf}\label{SF-E817-an-explicit-actual-singular-balanced-sheaf}

Use the CM product and integral divisor classes of SF-E815:

\[
 X=E_i^4,\qquad X_0=x+y,\qquad Y=\sqrt2(x-y),
 \qquad x^2=-\theta_+^2,\quad y^2=-\theta_-^2,
 \quad\int_Xx^2y^2=4.
\]

The rationality and integrality of \(X_0,Y\) come from the explicit
Hermitian matrices \(S\) and \(US\) in that theorem. Form the actual
polystable bundle

\[
 H=\bigl(L_{2X_0}\oplus L_{-2X_0}\bigr)^{\oplus3}
       \oplus\bigl(L_Y\oplus L_{-Y}\bigr)^{\oplus6}.
 \tag{E817.23}
\]

Its rank is eighteen and its odd character components vanish. Direct
expansion gives

\[
 \begin{aligned}
 \operatorname{ch}_2(H)
   &=12X_0^2+6Y^2=24(x^2+y^2),\\
 \operatorname{ch}_4(H)
   &=4X_0^4+\tfrac12Y^4=36x^2y^2.
 \end{aligned}
 \tag{E817.24}
\]

Choose sixteen distinct points \(Z\) and a nonzero linear functional on
the fibre of \(H\) at each point. These define an actual surjection
\(H\twoheadrightarrow\mathcal O_Z\). Let \(E\) be its kernel. It is
torsion-free of rank eighteen, and its double dual is \(H\). For
example, dualizing the sequence twice proves this directly because
\(\operatorname{Hom}(T,\mathcal O_X)=\operatorname{Ext}^1(T,\mathcal O_X)=0\).

Locally at a point of \(Z\), an invertible change of basis makes the
surjection \(R^{18}\to k\) the reduction of the first coordinate, so

\[
 E_p\simeq\mathfrak m_p\oplus R^{17}.
\]

The maximal ideal of a regular local ring of dimension four is not free
of rank one. Thus this is a genuinely singular torsion-free sheaf, not a
vector bundle.

Since \([\mathrm{pt}]=x^2y^2/4\), the entire character of its kernel is

\[
 \begin{split}
 \operatorname{ch}(E)
   &=18+24(x^2+y^2)+32x^2y^2\\
   &=18\left(1-\frac43\theta_+^2\right)
          \left(1-\frac43\theta_-^2\right)
     =18v_{8/3}.
 \end{split}
 \tag{E817.25}
\]

All the line summands of \(H\) have zero slope for the product
polarization, because their restrictions to each RM plane are primitive.
The bundle \(H\) is therefore polystable of slope zero. The kernel has
the same first Chern class, and every subsheaf of \(E\) is a subsheaf of
\(H\), so \(E\) is semistable of slope zero. Its singular stable factors
are not covered by the locally free Jordan--Hölder hypothesis in
SF-E813. The surjection can be chosen with all fibre components nonzero;
no direct-sum decomposition of the kernel is assumed.

Nevertheless SF-E817 excludes weak semiregularity of this actual sheaf
under every effective finite coherent action. In the ordinary category,
its sixteen support points give \(s=16\), hence the invariant-free full
trace kernel has dimension at least forty-nine by (E817.4). This lower
bound is not asserted to be sharp.

\subsection{Scope and source record}\label{SF-E817-scope-and-source-record}

This construction genuinely enters the singular-factor route and then
fails by the displayed exact maps. The universal theorem removes all
nonzero finite-length defects of locally free hulls in the balanced
character, without a stability or simplicity premise. It does not remove
defects supported on curves or surfaces, nor a reflexive hull that is
itself singular. Those cases have different local Ext groups and must
not be inferred from (E817.6).

The local vanishings are proved here by the four-variable Koszul complex
and finite-length filtrations. The only categorical input not reproved
is the already reread CY4 trace/evaluation adjunction and coherent
finite-action identity-sector assertion of supplied DEG\_L S-21. The
Fock identification and balanced rank eighteen are also rederived in the
displayed blocks and agree with the previously checked SF-E801. No new
preprint theorem, Hermitian--Einstein extension for singular sheaves, or
dg formality is used.
